% ── §2.3: Αναγνώριση Δραστηριότητας Ανθρώπου (HAR) ──────────

\section{Αναγνώριση Δραστηριότητας Ανθρώπου (HAR)}
\label{sec:har}

Η Αναγνώριση Δραστηριότητας Ανθρώπου (Human Activity Recognition, HAR)
αφορά την αυτόματη κατάταξη της φυσικής δραστηριότητας ενός ατόμου σε
προκαθορισμένες κλάσεις, βάσει δεδομένων αισθητήρων κινητής συσκευής
\cite{Anguita2013}. Στο πλαίσιο των συστημάτων AAL, το HAR επιτρέπει
την παρακολούθηση της λειτουργικής αυτονομίας χρηστών σε πραγματικό
χρόνο, χωρίς να απαιτείται φορετός εξοπλισμός ή υποδομή αισθητήρων
στο περιβάλλον \cite{Liu2021}.

\subsection{Αισθητήρες Κινητών Συσκευών}
\label{subsec:sensors}

Τα σύγχρονα smartphones διαθέτουν ενσωματωμένους αισθητήρες υψηλής
ευαισθησίας που καταγράφουν κινηματικές πληροφορίες σε πραγματικό χρόνο:

\begin{itemize}
    \item \textbf{Επιταχυνσιόμετρο (Accelerometer):} Μετρά την γραμμική
    επιτάχυνση κατά τους τρεις άξονες $\{x, y, z\}$ σε $\text{m/s}^2$.
    Συλλαμβάνει δυναμικά μοτίβα (βηματισμός, πτώση) αλλά και στατικές
    κατευθύνσεις (βαρύτητα ανάλογα με κλίση συσκευής).

    \item \textbf{Γυροσκόπιο (Gyroscope):} Μετρά τον ρυθμό γωνιακής
    στροφής κατά τους τρεις άξονες σε $\text{rad/s}$. Συμπληρώνει
    το επιταχυνσιόμετρο αποτυπώνοντας περιστροφικές κινήσεις που δεν
    ανιχνεύονται από τη γραμμική επιτάχυνση.
\end{itemize}

Ο συνδυασμός των δύο αισθητήρων παράγει ένα σήμα 6 καναλιών. Στο
UCI HAR dataset \cite{Anguita2013}, προστίθενται 3 ακόμα κανάλια
(επιταχύνσεις σώματος μετά από αφαίρεση βαρύτητας), οδηγώντας σε
είσοδο 9 καναλιών για τα αντίστοιχα μοντέλα μάθησης.

\subsection{Τεμαχισμός Σήματος με Συρόμενο Παράθυρο}
\label{subsec:sliding_window_bg}

Η ακατέργαστη χρονοσειρά τεμαχίζεται σε παράθυρα σταθερού μήκους με
τη μέθοδο \textit{sliding window}. Για μια χρονοσειρά $\mathbf{s}(t)$
δειγματοληψίας $f_s = 50\,\text{Hz}$, το παράθυρο $\mathbf{X}^{(k)}$
αρχίζει στο δείγμα $k \cdot (1-\rho) \cdot W$ και λήγει $W$ δείγματα
αργότερα, όπου $W$ το μήκος παραθύρου (σε δείγματα) και $\rho$ η
αναλογία επικάλυψης (overlap). Στο UCI HAR, $W = 128$ δείγματα ($2.56\,\text{s}$)
και $\rho = 0.5$, αποδίδοντας $\lfloor 2(T/W - 1) \rfloor + 1$ παράθυρα
ανά χρονοσειρά.

Κάθε παράθυρο $\mathbf{X}^{(k)} \in \mathbb{R}^{9 \times 128}$
αντιστοιχεί σε ένα δείγμα εισόδου για τον ταξινομητή. Αυτή ακριβώς
είναι η μορφή εισόδου $[B \times 9 \times 128]$ που δέχεται η
αρχιτεκτονική 1D-CNN (βλ.\ Ενότητα~\ref{sec:cnn_design}).

\subsection{Κλάσεις Δραστηριότητας}
\label{subsec:activity_classes}

Το UCI HAR dataset \cite{Anguita2013} περιλαμβάνει δεδομένα από 30
εθελοντές (ηλικία 19--48 ετών) για $K = 6$ κλάσεις δραστηριότητας:
\textit{περπάτημα} (Walking), \textit{ανάβαση σκαλιών} (Walking Upstairs),
\textit{κατάβαση σκαλιών} (Walking Downstairs), \textit{κάθισμα}
(Sitting), \textit{όρθια στάση} (Standing) και \textit{ξάπλωμα}
(Laying). Η κατηγοριοποίηση αυτή — τρεις δυναμικές και τρεις στατικές
κλάσεις — αποτελεί ευρέως αποδεκτό πρότυπο αξιολόγησης στη βιβλιογραφία
HAR \cite{Chen2021}.

\subsection{Προκλήσεις Αναγνώρισης σε Πραγματικές Συνθήκες}
\label{subsec:har_challenges}

Παρά τις επιδόσεις σε ελεγχόμενα πειραματικά περιβάλλοντα, το HAR
υπό πραγματικές συνθήκες AAL αντιμετωπίζει σημαντικές προκλήσεις
\cite{Hammerla2016}:

\begin{itemize}
    \item \textbf{Ατομική μεταβλητότητα:} Διαφορετικοί χρήστες
    εκτελούν ίδιες δραστηριότητες με εντελώς διαφορετικά κινηματικά
    μοτίβα — λόγω σωματότυπου, ηλικίας ή τρόπου κρατήματος της
    συσκευής. Αυτό οδηγεί σε Non-IID κατανομές δεδομένων μεταξύ
    clients στο FL σύστημα.

    \item \textbf{Θόρυβος αισθητήρων:} Ο θόρυβος αυξάνεται σε
    πραγματικά AAL σενάρια (άτυπη χρήση τηλεφώνου, κίνηση στο
    περιβάλλον) σε αντίθεση με τα πρωτόκολλα ελεγχόμενης συλλογής.

    \item \textbf{Ανισορροπία κλάσεων:} Στατικές δραστηριότητες
    (κάθισμα, ξάπλωμα) καταλαμβάνουν δυσανάλογα μεγαλύτερο χρόνο
    για ηλικιωμένους χρήστες, δημιουργώντας ανισορροπία στα τοπικά
    datasets των clients.
\end{itemize}

Οι ανωτέρω προκλήσεις επιβεβαιώνουν την ανάγκη για αξιολόγηση σε
πραγματικά δεδομένα — όπως το AAL dataset \cite{AAL_UCI2016}
— πέρα από τα ελεγχόμενα benchmarks. Ένας πρόσθετος παράγοντας
που επηρεάζει ειδικά τα FL συστήματα HAR είναι η ανάπτυξη σε
συσκευές περιορισμένων πόρων (smartphones, wearables): η υπολογιστική
ισχύς και η μνήμη είναι περιορισμένες, απαιτώντας αρχιτεκτονικές
μοντέλων που ισορροπούν ακρίβεια και αποδοτικότητα. Η εργασία των
Yao et al.\ \cite{Yao2017} (DeepIoT) δείχνει ότι η συμπίεση μοντέλων
για IoT συσκευές είναι εφικτή χωρίς σημαντική απώλεια ακρίβειας —
μια κατεύθυνση συμπληρωματική με αυτή της παρούσας εργασίας, η
οποία επιλέγει ήδη αρχιτεκτονική με περιορισμένο αριθμό παραμέτρων
($\approx 47$k) για τοπική εκπαίδευση. Παράλληλα, η χρήση πολλαπλών
τύπων αισθητήρων (fusion) για βελτίωση της αναγνώρισης σε ρεαλιστικά
σενάρια εξετάζεται από τους Gu et al.\ \cite{Gu2022}, αναδεικνύοντας
ότι η επιλογή των καναλιών εισόδου επηρεάζει σημαντικά την απόδοση
σε Non-IID κατανομές.
